Automated defect detection is a critical component of quality assurance in medical device manufacturing. Its importance has grown as medical devices have become increasingly complex and reliant on semiconductor components. As a consequence, \acrlong*{ml} has emerged as a promising approach for accelerating defect detection while improving the efficiency of quality assurance processes.

The effectiveness of \acrlong*{ml} methods depends heavily on the availability and quality of data, motivating researchers and organizations to focus on collecting large-scale datasets and organizing them in structured forms. In this context, \acrlongpl*{kg} have gained prominence as an effective approach to structuring and representing data, enabling efficient querying and facilitating the integration of information from diverse sources.

While aleatoric uncertainty arising from the data collection process has received increasing attention in recent studies, epistemic uncertainty associated with the decision model remains comparatively underexplored in the context of defect detection in medical device manufacturing. Detection errors can have serious consequences, ranging from potential harm during deployment to substantial costs associated with incorrectly identifying non-defective instances as defective. 

Motivated by these risks, we examine whether model reliability can be improved by abstaining from predictions associated with high epistemic uncertainty, using a simple yet principled rejection strategy. Finally, we evaluate the effectiveness of this approach on established synthetic benchmarks from the literature and real-world medical device report data.
The code to reproduce and extend our results is available at \url{https://code-repo.com}.

\raham{The clean repo must be linked here.}
